% GENERATED FILE. Edit canonical lessons, then run npm run generate:books.
% canonical-source-hash: fnv1a64:2f4f99bd632053db
% canonical-lessons: TA-C442-patuka, TA-C442-alivu, TA-C442-pacumai, TA-C442-cul, TA-C442-ucalatu

\chapter{Keeping It Green}
\label{ch:ta-keeping-it-green}
\hlchaptermodality{442}

\begin{chapteropening}
\textbf{By the end of this chapter:} \emph{I can talk about protecting nature, destruction and greenery, and floods surrounding a village.}

Complete the last lesson of chapter 442: I can talk about protecting nature, destruction and greenery, and floods surrounding a village.
\end{chapteropening}

\section[\texorpdfstring{\ta{பாதுகா} (pātukā)}{பாதுகா (pātukā)}]{\ta{பாதுகா} (pātukā) — to protect; to keep safe}

\label{lesson:TA-C442-patuka}



\begin{quote}
Before the new one: say the Tamil for to recycle, then the Tamil for a living creature.
\end{quote}

\subsection*{You'll want to know: \ta{பாதுகா}}
\textbf{\ta{பாதுகா}} — \emph{pātukā} — \textquotedblleft{}to protect; to keep safe\textquotedblright{}.

\begin{quote}
\ta{காடுகளைப்} \ta{பாதுகாக்க} \ta{அரசு} \ta{புதிய} \ta{சட்டம்} \ta{கொண்டு} \ta{வந்தது}.
\end{quote}

\begin{quote}
\emph{kāṭukaḷaip pātukākka aracu putiya caṭṭam koṇṭu vantatu.}
\end{quote}

\begin{quote}
The government brought a new law to protect the forests.
\end{quote}

Past (I …): \ta{பாதுகாத்தேன்}, \emph{pātukāttēṉ}.

\ta{பாதுகாப்பு}, safety, comes from it.

\begin{cousinweb}
\textbf{\ta{பாதுகா}} (\emph{pātukā}) ends in \ta{கா}, \textquotedblleft{}to guard, to protect\textquotedblright{} (DEDR 1416); the origin of its first part is not settled.
\end{cousinweb}

\begin{grammarlens}[title={What you've built}]
One more word for getting around.
\end{grammarlens}

\subsection*{Your turn}
\begin{itemize}
\raggedright
  \item \emph{Say it:} \emph{pātukā}
  \item \emph{Say it:} \emph{pātukā}, once more
  \item \emph{Say it:} \emph{kāṭukaḷaip pātukākka aracu putiya caṭṭam koṇṭu vantatu.} — The government brought a new law to protect the forests.
  \item \emph{Recall:} say the Tamil for to recycle, then the Tamil for a living creature, then say \emph{pātukā} again
\end{itemize}

\begin{tcolorbox}[breakable,colback=teal!4,colframe=teal!35!black,title={Before you move on}]
What does \ta{பாதுகா} mean? (\textquotedblleft{}To protect; to keep safe\textquotedblright{}.) Say it once more.
\end{tcolorbox}
\section[\texorpdfstring{\ta{அழிவு} (aḻivu)}{அழிவு (aḻivu)}]{\ta{அழிவு} (aḻivu) — destruction; ruin}

\label{lesson:TA-C442-alivu}



\begin{quote}
Before the new one: say the Tamil for a living creature, then the Tamil for to protect.
\end{quote}

\subsection*{You'll want to know: \ta{அழிவு}}
\textbf{\ta{அழிவு}} — \emph{aḻivu} — \textquotedblleft{}destruction; ruin\textquotedblright{}.

\begin{quote}
\ta{புயல்} \ta{கிராமத்தில்} \ta{பெரிய} \ta{அழிவை} \ta{ஏற்படுத்தியது}.
\end{quote}

\begin{quote}
\emph{puyal kirāmattil periya aḻivai ēṟpaṭuttiyatu.}
\end{quote}

\begin{quote}
The storm caused great destruction in the village.
\end{quote}

From \ta{அழி}, to perish, + -\ta{வு}.

\begin{cousinweb}
\textbf{\ta{அழிவு}} (\emph{aḻivu}) is native Dravidian (DEDR 277), the noun of \ta{அழி}, \textquotedblleft{}to perish, to be ruined\textquotedblright{}.
\end{cousinweb}

\begin{grammarlens}[title={What you've built}]
One more word for getting around.
\end{grammarlens}

\subsection*{Your turn}
\begin{itemize}
\raggedright
  \item \emph{Say it:} \emph{aḻivu}
  \item \emph{Say it:} \emph{aḻivu}, once more
  \item \emph{Say it:} \emph{puyal kirāmattil periya aḻivai ēṟpaṭuttiyatu.} — The storm caused great destruction in the village.
  \item \emph{Recall:} say the Tamil for a living creature, then the Tamil for to protect, then say \emph{aḻivu} again
\end{itemize}

\begin{tcolorbox}[breakable,colback=teal!4,colframe=teal!35!black,title={Before you move on}]
What does \ta{அழிவு} mean? (\textquotedblleft{}Destruction; ruin\textquotedblright{}.) Say it once more.
\end{tcolorbox}
\section[\texorpdfstring{\ta{பசுமை} (pacumai)}{பசுமை (pacumai)}]{\ta{பசுமை} (pacumai) — greenery; freshness}

\label{lesson:TA-C442-pacumai}



\begin{quote}
Before the new one: say the Tamil for to protect, then the Tamil for destruction.
\end{quote}

\subsection*{You'll want to know: \ta{பசுமை}}
\textbf{\ta{பசுமை}} — \emph{pacumai} — \textquotedblleft{}greenery; freshness\textquotedblright{}.

\begin{quote}
\ta{மழைக்குப்} \ta{பிறகு} \ta{மலைகள்} \ta{பசுமையாக} \ta{இருந்தன}.
\end{quote}

\begin{quote}
\emph{maḻaikkup piṟaku malaikaḷ pacumaiyāka iruntaṉa.}
\end{quote}

\begin{quote}
After the rain the hills were green.
\end{quote}

With -\ta{ஆக} it describes: \ta{பசுமையாக}, green and fresh.

\begin{cousinweb}
\textbf{\ta{பசுமை}} (\emph{pacumai}) is native Dravidian (DEDR 3821), \textquotedblleft{}greenness, verdure, freshness\textquotedblright{}.
\end{cousinweb}

\begin{grammarlens}[title={What you've built}]
One more word for getting around.
\end{grammarlens}

\subsection*{Your turn}
\begin{itemize}
\raggedright
  \item \emph{Say it:} \emph{pacumai}
  \item \emph{Say it:} \emph{pacumai}, once more
  \item \emph{Say it:} \emph{maḻaikkup piṟaku malaikaḷ pacumaiyāka iruntaṉa.} — After the rain the hills were green.
  \item \emph{Recall:} say the Tamil for to protect, then the Tamil for destruction, then say \emph{pacumai} again
\end{itemize}

\begin{tcolorbox}[breakable,colback=teal!4,colframe=teal!35!black,title={Before you move on}]
What does \ta{பசுமை} mean? (\textquotedblleft{}Greenery; freshness\textquotedblright{}.) Say it once more.
\end{tcolorbox}
\section[\texorpdfstring{\ta{சூழ்} (cūḻ)}{சூழ் (cūḻ)}]{\ta{சூழ்} (cūḻ) — to surround; to envelop}

\label{lesson:TA-C442-cul}



\begin{quote}
Before the new one: say the Tamil for destruction, then the Tamil for greenery.
\end{quote}

\subsection*{You'll want to know: \ta{சூழ்}}
\textbf{\ta{சூழ்}} — \emph{cūḻ} — \textquotedblleft{}to surround; to envelop\textquotedblright{}.

\begin{quote}
\ta{வெள்ளம்} \ta{கிராமத்தைச்} \ta{சூழ்ந்தது}.
\end{quote}

\begin{quote}
\emph{veḷḷam kirāmattaic cūḻntatu.}
\end{quote}

\begin{quote}
The flood surrounded the village.
\end{quote}

Past (it …): \ta{சூழ்ந்தது}, \emph{cūḻntatu}.

What is surrounded takes -\ta{ஐ}: \ta{கிராமத்தைச்} \ta{சூழ்}.

\begin{cousinweb}
\textbf{\ta{சூழ்}} (\emph{cūḻ}) is native Dravidian (DEDR 2698), \textquotedblleft{}to encompass, to surround, to envelop\textquotedblright{}.
\end{cousinweb}

\begin{grammarlens}[title={What you've built}]
One more word for getting around.
\end{grammarlens}

\subsection*{Your turn}
\begin{itemize}
\raggedright
  \item \emph{Say it:} \emph{cūḻ}
  \item \emph{Say it:} \emph{cūḻ}, once more
  \item \emph{Say it:} \emph{veḷḷam kirāmattaic cūḻntatu.} — The flood surrounded the village.
  \item \emph{Recall:} say the Tamil for destruction, then the Tamil for greenery, then say \emph{cūḻ} again
\end{itemize}

\begin{tcolorbox}[breakable,colback=teal!4,colframe=teal!35!black,title={Before you move on}]
What does \ta{சூழ்} mean? (\textquotedblleft{}To surround; to envelop\textquotedblright{}.) Say it once more.
\end{tcolorbox}
\section[\texorpdfstring{\ta{ஊசலாடு} (ūcalāṭu)}{ஊசலாடு (ūcalāṭu)}]{\ta{ஊசலாடு} (ūcalāṭu) — to swing about; to dangle}

\label{lesson:TA-C442-ucalatu}



\begin{quote}
Before the new one: say the Tamil for greenery, then the Tamil for to surround.
\end{quote}

\subsection*{You'll want to know: \ta{ஊசலாடு}}
\textbf{\ta{ஊசலாடு}} — \emph{ūcalāṭu} — \textquotedblleft{}to swing about; to dangle\textquotedblright{}.

\begin{quote}
\ta{காற்றில்} \ta{கயிற்றுப்} \ta{பாலம்} \ta{ஊசலாடியது}.
\end{quote}

\begin{quote}
\emph{kāṟṟil kayiṟṟup pālam ūcalāṭiyatu.}
\end{quote}

\begin{quote}
The rope bridge swung about in the wind.
\end{quote}

Past (it …): \ta{ஊசலாடியது}, \emph{ūcalāṭiyatu}.

\ta{ஊசல்}, a swing, + \ta{ஆடு}, to move.

\begin{cousinweb}
\textbf{\ta{ஊசலாடு}} (\emph{ūcalāṭu}) is built from \ta{ஊசல்}, \textquotedblleft{}moving to and fro, a swing\textquotedblright{} (DEDR 731), and \ta{ஆடு}, \textquotedblleft{}to move, to swing\textquotedblright{} (DEDR 347).
\end{cousinweb}

\begin{grammarlens}[title={What you've built}]
One more word for getting around.
\end{grammarlens}

\subsection*{Your turn}
\begin{itemize}
\raggedright
  \item \emph{Say it:} \emph{ūcalāṭu}
  \item \emph{Say it:} \emph{ūcalāṭu}, once more
  \item \emph{Say it:} \emph{kāṟṟil kayiṟṟup pālam ūcalāṭiyatu.} — The rope bridge swung about in the wind.
  \item \emph{Recall:} say the Tamil for greenery, then the Tamil for to surround, then say \emph{ūcalāṭu} again
\end{itemize}

\begin{tcolorbox}[breakable,colback=teal!4,colframe=teal!35!black,title={Before you move on}]
What does \ta{ஊசலாடு} mean? (\textquotedblleft{}To swing about; to dangle\textquotedblright{}.) Say it once more.
\end{tcolorbox}
